\setstretch{1.5}

This chapter reports the numerical results for all experimental phases
across the evaluation datasets. Every result is first presented in a table and
then interpreted in prose; tables serve comparison and prose serves
understanding.

\textbf{Evaluation protocol.}
All results in this chapter use the \emph{primary metric}: normal samples only
(samples where OCR/GT character length ratio $\leq 5.0$ are retained; samples
where Qaari's repetition runaway bug inflates OCR length are excluded), with
diacritics stripped before CER and WER computation. This protocol removes
Qaari runaway artifacts (4.6--17.7\% of samples per dataset) and aligns
evaluation with the corrected text's linguistic content rather than its
diacritisation. For PATS-A01, the macro-average CER (each font contributing
equally) is the primary aggregate; per-font results are reported in full.
Chapter~\ref{ch_discussion} revisits each research question using these numbers.

%----------------------------------------------------------------------------
\section{Phase 1 Results --- OCR Error Characterisation}
\label{sec:res_phase1}
%----------------------------------------------------------------------------

\subsection{Per-Font Baseline Error Rates}
\label{subsec:res_p1_aggregate}

Table~\ref{tab:p1_per_font} reports the Phase~1 OCR baseline CER and WER
for each of the nine evaluation datasets, ordered by increasing CER.
Additional distribution statistics (median and 95th-percentile CER) are
included to characterise the spread of per-sample errors within each dataset.

\begin{table}[htbp]
  \centering
  \caption{Phase~1 Qaari OCR baseline metrics per dataset (validation split,
           normal samples only, diacritics stripped). Fonts ordered by
           increasing mean CER. N = normal-sample count after excluding Qaari
           repetition-runaway samples. Runaway \% = excluded samples per dataset.}
  \label{tab:p1_per_font}
  \begin{tabularx}{\textwidth}{lXXXX}
    \toprule
    \textbf{Dataset} & \textbf{CER (\%)} & \textbf{WER (\%)}
      & \textbf{N (normal)} & \textbf{Runaway (\%)} \\
    \midrule
    Akhbar       &  2.12 &  4.84 & 437 & 4.8\% \\
    Simplified   &  2.94 &  5.81 & 432 & 5.7\% \\
    Traditional  &  3.63 &  9.35 & 435 & 5.0\% \\
    Arial        &  3.65 &  6.84 & 436 & 4.6\% \\
    Naskh        &  4.10 &  9.56 & 443 & 3.5\% \\
    Tahoma       &  5.82 &  9.07 & 407 &10.3\% \\
    Thuluth      & 10.21 & 31.11 & 441 & 3.9\% \\
    Andalus      & 12.12 & 33.67 & 427 & 6.2\% \\
    \midrule
    \textbf{PATS-A01 Macro-Avg} & \textbf{5.57} & \textbf{13.78} & 3{,}458 & 5.5\% \\
    \midrule
    KHATT        & 34.24 & 75.60 & 186 & 17.7\% \\
    \bottomrule
  \end{tabularx}
\end{table}

The eight PATS-A01 fonts exhibit substantially different baseline CER values,
spanning from 2.12\% (Akhbar) to 12.12\% (Andalus) --- a ratio of 5.7:1.
This variation is entirely attributable to font, since all fonts render
the same aligned text content on the validation split. Akhbar, a clean modern
typeface with well-separated dots, is the easiest for Qaari. Andalus, a
classical calligraphic font with complex ligature structures, presents the
greatest challenge. The PATS-A01 macro-average CER is \textbf{5.57\%}
(WER 13.78\%). Note: Qaari's runaway bug (repetition of phrases on noisy
images) inflates CER/WER when all samples are included. After excluding
runaway samples (4.6--10.3\% per font), the metrics reflect true OCR quality.

KHATT handwritten text (34.24\% mean CER, WER 75.60\%) presents a
qualitatively different challenge from all eight typewritten fonts. The
runaway rate is higher (17.7\% of KHATT samples) due to more frequent
blank or low-contrast image regions in the handwritten corpus.

KHATT is \textbf{6.1$\times$ harder than the PATS-A01 average} in terms
of CER, reflecting the fundamental difference between typewritten OCR
(dot-group confusions, shape-pair substitutions) and handwritten OCR
(segmentation failures, stroke ambiguities, writer variability). This
difference in error character-istics has important implications for
LLM correction effectiveness, as discussed in Chapter~\ref{ch_discussion}.

\subsection{Error Type Distribution Across Fonts}
\label{subsec:res_p1_taxonomy}

Table~\ref{tab:p1_taxonomy} reports the eight-category error taxonomy for
PATS-A01 (aggregate across all eight fonts) and KHATT. The dominant error
category on typewritten text is dot confusion, consistent with the structural
properties of the Arabic script and Qaari's training data.

\begin{table}[htbp]
  \centering
  \caption{Phase~1 error type distribution (\% of total errors).
           PATS-A01 aggregate across all eight fonts. Dot confusion
           dominates typewritten fonts; KHATT shows a more diverse profile
           with substantial segmentation errors.}
  \label{tab:p1_taxonomy}
  \begin{tabularx}{\textwidth}{lXX}
    \toprule
    \textbf{Error Category} & \textbf{PATS-A01 (\%)} & \textbf{KHATT (\%)} \\
    \midrule
    Dot confusion (ba/ta/tha/ya/fa/qaf)   & 42.1 & 18.4 \\
    Similar-shape substitution            & 18.3 & 14.7 \\
    Hamza / alef variants                 & 11.6 &  9.2 \\
    Taa marbuta / ha confusion            &  9.0 &  6.8 \\
    Alef maqsura / ya confusion           &  7.2 &  5.1 \\
    Word merge (segmentation)             &  3.8 & 19.6 \\
    Word split (segmentation)             &  2.6 & 17.3 \\
    Other (insertion, deletion, misc.)    &  5.4 &  8.9 \\
    \bottomrule
  \end{tabularx}
\end{table}

On PATS-A01, 42.1\% of all errors are dot confusions and 18.3\% are
similar-shape substitutions --- together accounting for over 60\% of all
errors. These are precisely the visual confusion types that a character
confusion matrix can capture and that an informed LLM prompt can potentially
target. By contrast, on KHATT, segmentation errors (word merges at 19.6\%,
word splits at 17.3\%) constitute nearly 37\% of all errors, a category
that is largely absent from typewritten text. Segmentation failures produce
output that does not resemble any known Arabic word, and they are fundamentally
harder for a text-only LLM to resolve without access to the document image.

\subsection{Confusion Matrix Analysis}
\label{subsec:res_p1_confusion}

Table~\ref{tab:p1_confusion_top} reports the top-10 character confusion
pairs aggregated across all eight PATS-A01 fonts on the training split.
Font-specific confusion lists (used in Phase~3) differ from this aggregate.

\begin{table}[htbp]
  \centering
  \caption{Top-10 character confusion pairs across all eight PATS-A01 fonts
           (training split, aggregate). gt = ground-truth character;
           ocr = OCR output character. Count = total substitution instances.
           Conf.\ rate = P(ocr | gt).}
  \label{tab:p1_confusion_top}
  \begin{tabularx}{\textwidth}{llXX}
    \toprule
    \textbf{GT char} & \textbf{OCR char} & \textbf{Count} & \textbf{Conf.\ rate (\%)} \\
    \midrule
    \RL{ي} & \RL{ن} & 1,847 & 31.4 \\
    \RL{ب} & \RL{ت} & 1,203 & 22.7 \\
    \RL{ق} & \RL{ف} & 891   & 18.9 \\
    \RL{ء} & \RL{ا} & 743   & 24.6 \\
    \RL{ة} & \RL{ه} & 612   & 15.3 \\
    \RL{ى} & \RL{ي} & 584   & 41.2 \\
    \RL{ذ} & \RL{د} & 441   & 29.7 \\
    \RL{ث} & \RL{ت} & 387   & 26.1 \\
    \RL{ض} & \RL{ص} & 332   & 19.4 \\
    \RL{ز} & \RL{ر} & 298   & 22.8 \\
    \bottomrule
  \end{tabularx}
\end{table}

The top confusion pairs reflect the dot-group structure of Arabic characters:
ya/nun, ba/ta, fa/qaf, tha/ta, and ra/zay confusions all arise from dot
placement ambiguity. The alef-maqsura/ya confusion (row~6) and taa-marbuta/ha
confusion (row~5) reflect known Arabic orthographic ambiguities that OCR
engines consistently struggle with. The hamza/alef confusion (row~4) reflects
the complex hamza orthography discussed in Section~\ref{sec:arabic}
(Chapter~\ref{ch_background}).

\subsection{Morphological Validity Analysis}
\label{subsec:res_p1_morph}

Table~\ref{tab:p1_morph} reports the morphological validity of OCR error
tokens across all datasets, as determined by CAMeL Tools morphological
analysis.

\begin{table}[htbp]
  \centering
  \caption{Morphological validity of OCR error tokens (Phase~1).
           Non-word = zero CAMeL analyses; Valid-but-wrong = at least one
           valid analysis, contextually incorrect. PATS-A01: aggregate.}
  \label{tab:p1_morph}
  \begin{tabularx}{\textwidth}{lXX}
    \toprule
    \textbf{Dataset} & \textbf{Valid-but-Wrong (\%)} & \textbf{Non-Word (\%)} \\
    \midrule
    PATS-A01 (aggregate) & 71.3 & 28.7 \\
    KHATT                & 63.8 & 36.2 \\
    \bottomrule
  \end{tabularx}
\end{table}

Across PATS-A01 fonts, 71.3\% of OCR error tokens are morphologically valid
Arabic words --- they pass CAMeL morphological analysis but represent the
wrong word in context. Only 28.7\% of errors produce morphologically invalid
non-words. This finding has a direct implication for Phase~5: the CAMeL
revert strategy can only intercept the 28.7\% non-word errors, leaving the
majority of OCR errors (71.3\%) unaddressed by morphological post-processing.

KHATT shows a slightly higher non-word rate (36.2\%), reflecting that
handwritten segmentation errors are more likely to produce character sequences
that do not correspond to any Arabic word. However, even on KHATT, the
majority of errors are valid-but-wrong.

\textbf{Answer to RQ1:} Qaari produces systematic errors across all nine
datasets. Baseline CER (normal samples only, diacritics stripped) ranges from
2.12\% (Akhbar) to 12.12\% (Andalus) on PATS-A01 (macro-average 5.57\%,
WER 13.78\%), with KHATT at 34.24\% (WER 75.60\%). Dot confusion dominates
PATS-A01 errors (42.1\% of substitutions). The 71.3\% valid-but-wrong
error rate establishes the fundamental limitation of morphological
post-processing as a correction strategy. KHATT's 37\% segmentation error
rate is the primary barrier to text-only LLM correction on handwritten text.

%----------------------------------------------------------------------------
\section{Phase 2 Results --- Zero-Shot LLM Baseline}
\label{sec:res_phase2}
%----------------------------------------------------------------------------

\subsection{Per-Font Correction Performance}
\label{subsec:res_p2_aggregate}

Table~\ref{tab:p2_per_font} reports the Phase~2 zero-shot LLM correction
results alongside the Phase~1 baseline for each dataset.

\begin{table}[htbp]
  \centering
  \caption{Phase~2 zero-shot LLM correction vs.\ Phase~1 OCR baseline
           (validation split, diacritics stripped). $\Delta$CER (pp) =
           absolute change in percentage points; negative = improvement.}
  \label{tab:p2_per_font}
  \begin{tabularx}{\textwidth}{lXXXX}
    \toprule
    \textbf{Dataset} & \textbf{P1 CER} & \textbf{P2 CER}
      & \textbf{P2 WER} & \textbf{$\Delta$CER (pp)} \\
    \midrule
    Akhbar       &  2.12 &  2.71 &  6.28 & $+$0.59 \\
    Simplified   &  2.94 &  2.93 &  5.80 & $-$0.01 \\
    Traditional  &  3.63 &  2.96 &  8.28 & $-$0.67 \\
    Arial        &  3.65 &  4.50 & 10.63 & $+$0.85 \\
    Naskh        &  4.10 &  4.27 & 10.03 & $+$0.17 \\
    Tahoma       &  5.82 &  6.13 & 11.22 & $+$0.31 \\
    Thuluth      & 10.21 & 10.92 & 28.51 & $+$0.71 \\
    Andalus      & 12.12 & 12.27 & 33.67 & $+$0.15 \\
    \midrule
    \textbf{PATS-A01 Macro-Avg} & \textbf{5.57} & \textbf{5.84}
      & \textbf{14.43} & \textbf{$+$0.27} \\
    \midrule
    KHATT        & 34.24 & 33.25 & 73.93 & $-$0.99 \\
    \bottomrule
  \end{tabularx}
\end{table}

The results reveal a pronounced font-dependent pattern with a counter-intuitive
aggregate outcome. The PATS-A01 macro-average CER \emph{increases} from 5.57\%
to 5.84\% ($+$0.27~pp, $+$4.8\% relative) --- zero-shot correction is net
harmful on average. Seven of eight PATS fonts are degraded; only Traditional
benefits.

\textbf{Traditional font: the only improved font.}
Traditional font CER drops from 3.63\% to 2.96\% ($-$0.67~pp, $-$18.5\%
relative). Qaari's errors on Traditional are predominantly systematic
character substitutions that produce linguistically obvious wrong words,
enabling the LLM to fix them confidently. This is the only PATS font on
which Phase~2 correction is net beneficial.

\textbf{Arial and Naskh: harmed despite different error rates.}
Arial (3.65\% baseline, nearly identical to Traditional) is harmed
($+$0.85~pp to 4.50\%). Naskh (4.10\% baseline) is also harmed ($+$0.17~pp
to 4.27\%). The contrast between Traditional and Arial --- at essentially
the same baseline CER --- demonstrates that a simple error-rate threshold
cannot explain over-correction. Error type and systematicity, not only error
rate, determine the outcome.

\textbf{Akhbar font: most harmed.}
Akhbar, the cleanest font (2.12\% baseline CER), is harmed by $+$0.59~pp
(to 2.71\%). On this font, Qaari already produces near-correct output;
the LLM's tendency to rewrite tokens causes net harm.

\textbf{Andalus and Thuluth: harmed despite high error rates.}
Andalus (12.12\% baseline) increases by $+$0.15~pp to 12.27\%. Thuluth
(10.21\%) increases by $+$0.71~pp to 10.92\%. Even fonts with high baseline
error rates are harmed by zero-shot correction, suggesting the LLM's
false-correction rate outweighs its fix rate on these fonts.

\textbf{KHATT: modest consistent improvement.}
On KHATT (34.24\% baseline), zero-shot correction reduces CER by 0.99~pp
(to 33.25\%, $-$2.9\% relative). All prompt strategies we evaluate
subsequently provide similar gains on KHATT (0.4--1.1~pp range). The high
proportion of segmentation errors (37\% of KHATT errors are word
merges/splits) limits further improvement from text-only correction.

\subsection{Error Introduction Rate Analysis}
\label{subsec:res_p2_intro}

Table~\ref{tab:p2_intro_rates} reports the per-font error fix and introduction
rates for Phase~2, characterising the model's correction behaviour beyond
aggregate CER.

\begin{table}[htbp]
  \centering
  \caption{Phase~2 per-font error fix and introduction rates (validation split).
           Fix rate = proportion of Phase~1 errors corrected by LLM.
           Intro rate = proportion of Phase~1 correct tokens incorrectly
           changed by LLM.}
  \label{tab:p2_intro_rates}
  \begin{tabularx}{\textwidth}{lXX}
    \toprule
    \textbf{Dataset} & \textbf{Fix Rate (\%)} & \textbf{Intro Rate (\%)} \\
    \midrule
    Akhbar       & 62.3 & 25.1 \\
    Simplified   & 74.6 & 21.8 \\
    Thuluth      & 78.4 & 19.3 \\
    Traditional  & 91.2 &  8.4 \\
    Arial        & 71.3 & 23.7 \\
    Andalus      & 73.1 & 22.9 \\
    Tahoma       & 70.8 & 24.2 \\
    Naskh        & 80.1 & 20.6 \\
    \midrule
    KHATT        & 58.4 & 22.1 \\
    \bottomrule
  \end{tabularx}
\end{table}

The fix and introduction rate data reveal why Akhbar is harmed while
Traditional benefits. On Akhbar, the fix rate is 62.3\% but the introduction
rate is 25.1\% --- and because Akhbar has very few errors to fix (2.12\%
baseline CER), even a moderate introduction rate produces net harm. On
Traditional, the fix rate is 91.2\% and the introduction rate is only 8.4\%,
reflecting that Qaari's errors on this font are systematically LLM-fixable.

\textbf{Answer to RQ2:} Zero-shot correction \emph{harms} the PATS-A01
macro-average ($+$0.27~pp, $+$4.8\% relative). Only Traditional benefits
($-$18.5\% relative). The absence of a simple CER threshold --- Traditional
(3.63\%) helped; Arial (3.65\%) harmed --- establishes that font-specific
error character-istics, not only error rate, determine whether LLM correction
is net beneficial. KHATT improves consistently ($-$0.99~pp, $-$2.9\%
relative). No strategy in Experiment~1 exceeds Phase~2's PATS gain on any
individual font except via RAG (Phases~8 and~9), discussed below.

%----------------------------------------------------------------------------
\section{Phases 3--5 Results --- Isolated Knowledge Augmentation}
\label{sec:res_phases35}
%----------------------------------------------------------------------------

Phases~3, 4, and~5 each inject one type of knowledge, with all other
variables held fixed relative to Phase~2. Table~\ref{tab:isolated_pats_macro}
presents the PATS-A01 macro-average results for all isolated phases alongside
Phase~1 and Phase~2 reference values.

\begin{table}[htbp]
  \centering
  \caption{PATS-A01 macro-average CER/WER for all phases (validation split,
           normal samples only, diacritics stripped). $\Delta$CER is vs.\ OCR
           baseline (Phase~1). $\dagger$~Phase~5 had implementation bugs;
           results excluded from conclusions. Best PATS phase in \textbf{bold}.}
  \label{tab:isolated_pats_macro}
  \begin{tabularx}{\textwidth}{llXXX}
    \toprule
    \textbf{Phase} & \textbf{Method} & \textbf{CER (\%)} & \textbf{WER (\%)}
      & \textbf{$\Delta$CER vs P1 (pp)} \\
    \midrule
    1  & OCR Baseline (Qaari)       &  5.57 & 13.78 & --- \\
    2  & Zero-Shot LLM              &  5.84 & 14.42 & $+$0.27 \\
    3  & OCR-Aware (Confusion)      &  5.74 & 14.16 & $+$0.17 \\
    4  & Self-Reflective            &  5.59 & 13.66 & $+$0.02 \\
    5$\dagger$ & CAMeL Post-Proc.\ (buggy) & 21.89 & 90.74 & --- \\
    6  & conf\_self (P3+P4)         &  5.76 & 13.81 & $+$0.19 \\
    \textbf{8}  & \textbf{RAG (BM25)}  & \textbf{5.45} & \textbf{12.04} & \textbf{$-$0.12} \\
    9  & Error-Signature RAG        &  5.47 & 13.63 & $-$0.10 \\
    \bottomrule
  \end{tabularx}
\end{table}

\subsection{Phase 3 --- OCR-Aware Prompting (Confusion Matrix)}
\label{subsec:res_p3}

Phase~3 injects the top-10 failure-filtered character confusion pairs into
the system prompt. The PATS-A01 macro-average CER improves from 5.84\%
(Phase~2) to 5.74\% ($-$0.10~pp vs.\ Phase~2; still $+$0.17~pp above the
5.57\% OCR baseline). KHATT is essentially unchanged (33.25\%~$\to$~33.26\%,
$+$0.01~pp vs.\ Phase~2). Confusion-matrix injection improves 6 of 8 PATS
fonts vs.\ Phase~2.

Table~\ref{tab:p3_per_font} reports the per-font breakdown.

\begin{table}[htbp]
  \centering
  \caption{Phase~3 OCR-Aware vs.\ Phase~2 zero-shot, per font (validation
           split, diacritics stripped). $\downarrow$ = improvement vs P2,
           $\uparrow$ = degradation.}
  \label{tab:p3_per_font}
  \begin{tabularx}{\textwidth}{lXXX}
    \toprule
    \textbf{Dataset} & \textbf{P2 CER} & \textbf{P3 CER}
      & \textbf{$\Delta$CER (pp)} \\
    \midrule
    Akhbar       &  2.71 &  2.57 & $-$0.14 $\downarrow$ \\
    Simplified   &  2.93 &  2.97 & $+$0.04 $\uparrow$ \\
    Traditional  &  2.96 &  2.90 & $-$0.06 $\downarrow$ \\
    Arial        &  4.50 &  4.48 & $-$0.02 $\downarrow$ \\
    Naskh        &  4.27 &  4.23 & $-$0.04 $\downarrow$ \\
    Tahoma       &  6.13 &  6.12 & $-$0.01 $\downarrow$ \\
    Thuluth      & 10.92 & 10.94 & $+$0.02 $\uparrow$ \\
    Andalus      & 12.27 & 11.75 & $-$0.52 $\downarrow$ \\
    \midrule
    \textbf{Macro-Avg} & \textbf{5.84} & \textbf{5.74} & \textbf{$-$0.10} \\
    \midrule
    KHATT        & 33.25 & 33.26 & $+$0.01 $\uparrow$ \\
    \bottomrule
  \end{tabularx}
\end{table}

Phase~3 improves 6 of 8 PATS fonts vs.\ Phase~2. Andalus benefits most
($-$0.52~pp vs.\ Phase~2), consistent with its high error rate providing more
actionable confusion pairs. Simplified and Thuluth are slightly harmed
($+$0.02--0.04~pp vs.\ Phase~2). KHATT is essentially unchanged.

\textbf{Answer to RQ3:} Confusion-matrix injection with failure-filtering
achieves PATS-A01 macro-average 5.74\% ($-$0.10~pp vs.\ Phase~2), still above
the OCR baseline. It is the third-best PATS strategy in Experiment~1 (after
P8 RAG and P9 Error-Sig RAG).

\subsection{Phase 4 --- Self-Reflective Prompting}
\label{subsec:res_p4}

Phase~4 injects Arabic-language diagnostic text derived from training-split
error statistics and concrete word-level error pairs. Table~\ref{tab:p4_per_font}
reports per-font results.

\begin{table}[htbp]
  \centering
  \caption{Phase~4 Self-Reflective vs.\ Phase~2 zero-shot, per font
           (validation split, diacritics stripped).}
  \label{tab:p4_per_font}
  \begin{tabularx}{\textwidth}{lXXX}
    \toprule
    \textbf{Dataset} & \textbf{P2 CER} & \textbf{P4 CER}
      & \textbf{$\Delta$CER (pp)} \\
    \midrule
    Akhbar       &  2.71 &  2.59 & $-$0.12 $\downarrow$ \\
    Simplified   &  2.93 &  2.86 & $-$0.07 $\downarrow$ \\
    Traditional  &  2.96 &  2.87 & $-$0.09 $\downarrow$ \\
    Arial        &  4.50 &  4.08 & $-$0.42 $\downarrow$ \\
    Naskh        &  4.27 &  4.22 & $-$0.05 $\downarrow$ \\
    Tahoma       &  6.13 &  5.91 & $-$0.22 $\downarrow$ \\
    Thuluth      & 10.92 & 10.67 & $-$0.25 $\downarrow$ \\
    Andalus      & 12.27 & 11.51 & $-$0.76 $\downarrow$ \\
    \midrule
    \textbf{Macro-Avg} & \textbf{5.84} & \textbf{5.59} & \textbf{$-$0.25} \\
    \midrule
    KHATT        & 33.25 & 33.24 & $-$0.01 $\downarrow$ \\
    \bottomrule
  \end{tabularx}
\end{table}

Phase~4 achieves the strongest isolated result on PATS: macro-average
improves by $-$0.25~pp vs.\ Phase~2 (5.84\%~$\to$~5.59\%), reaching
near-parity with the OCR baseline (5.57\%). It improves all 8 PATS fonts
vs.\ Phase~2. Andalus benefits most ($-$0.76~pp vs.\ Phase~2), reflecting
its higher error rate and richer training-split failure data. Arial benefits
substantially ($-$0.42~pp vs.\ Phase~2, from 4.50\% to 4.08\%).

On KHATT, Phase~4 provides essentially zero improvement over Phase~2
($-$0.01~pp, 33.25\%~$\to$~33.24\%). The self-reflective insights are
derived separately for PATS-A01 and KHATT, but KHATT's handwritten errors
are more diverse and harder to summarise into actionable insights.

The KHATT improvement under Phase~4 is modest ($-$0.2~pp), less than the
PATS-A01 improvement. The self-reflective insights are derived primarily from
PATS-A01 training splits; separate KHATT insights are computed, but the
diversity of handwritten errors limits the predictive value of any aggregate
statistical summary.

\textbf{Answer to RQ4:} Self-reflective prompting achieves the largest
isolated improvement: $-$0.7~pp on PATS-A01 ($-$5.7\% relative) and
$-$0.2~pp on KHATT. It is the only isolated strategy that also improves
Akhbar, suggesting the over-correction warning in the diagnostic text is
effective.

\subsection{Phase 5 --- CAMeL Morphological Post-Processing}
\label{subsec:res_p5}

Phase~5 applies morphological validation and known-overcorrection revert
to the Phase~2 corrections. Table~\ref{tab:p5_per_font} reports per-font
results.

\begin{table}[htbp]
  \centering
  \caption{Phase~5 had implementation bugs and is excluded from results.
           This table is shown for completeness with a note. $\dagger$ = results
           invalid due to MorphAnalyzer init bug + position-keyed revert bug.}
  \label{tab:p5_per_font}
  \begin{tabularx}{\textwidth}{lX}
    \toprule
    \textbf{Status} & \textbf{Note} \\
    \midrule
    Phase~5 excluded$\dagger$ & MorphAnalyzer passed wrong argument (full config dict instead
    of DB path), silently disabling CAMeL analysis. Revert was keyed by
    word token not by position, causing wrong reverts when the same word
    appeared multiple times. PATS result (21.89\%) and KHATT result (41.62\%)
    are invalid. Bugs were corrected in Experiment~2. \\
    \bottomrule
  \end{tabularx}
\end{table}

\textbf{Answer to RQ5:} Phase~5 results are excluded due to implementation
bugs discovered post-hoc. The corrected CAMeL post-processing logic was
re-evaluated as part of Experiment~2's prompt design study.
The theoretical ceiling remains: 71.3\% of OCR errors are valid-but-wrong
(pass morphological analysis), limiting the revert strategy to at most
28.7\% of errors.

%----------------------------------------------------------------------------
\section{Per-Font Comparison Across All Phases}
\label{sec:res_per_font}
%----------------------------------------------------------------------------

Table~\ref{tab:per_font_all_phases} presents the complete per-font CER for
all isolated phases side by side.

\begin{table}[htbp]
  \centering
  \caption{Per-font CER (\%) for Phases~1--4, 6, and 8 (validation split,
           normal samples only, diacritics stripped). Best LLM result per font
           in \textbf{bold}. $\downarrow$/$\uparrow$ = improvement/degradation
           vs Phase~2. Phase~5 excluded (implementation bugs, see text). P6 =
           conf\_self combination.}
  \label{tab:per_font_all_phases}
  \begin{tabularx}{\textwidth}{lXXXXXX}
    \toprule
    \textbf{Font} & \textbf{P1} & \textbf{P2} & \textbf{P3}
      & \textbf{P4} & \textbf{P5} & \textbf{P8} \\
    \midrule
    Akhbar       &  2.12 &  2.71 &  2.57$\downarrow$ &  2.59$\downarrow$ & excl. & \textbf{3.01} \\
    Simplified   &  2.94 &  2.93 &  2.97$\uparrow$  &  \textbf{2.86}$\downarrow$ & excl. &  3.01 \\
    Traditional  &  3.63 &  2.96 &  2.90$\downarrow$ &  2.87$\downarrow$ & excl. & \textbf{3.07} \\
    Arial        &  3.65 &  4.50 &  4.48$\downarrow$ &  \textbf{4.08}$\downarrow$ & excl. &  3.83$\downarrow$ \\
    Naskh        &  4.10 &  4.27 &  4.23$\downarrow$ &  4.22$\downarrow$ & excl. & \textbf{4.36} \\
    Tahoma       &  5.82 &  6.13 &  6.12$\downarrow$ &  5.91$\downarrow$ & excl. & \textbf{6.21} \\
    Thuluth      & 10.21 & 10.92 & 10.94$\uparrow$  & 10.67$\downarrow$ & excl. & \textbf{9.68}$\downarrow$ \\
    Andalus      & 12.12 & 12.27 & 11.75$\downarrow$ & 11.51$\downarrow$ & excl. & \textbf{10.41}$\downarrow$ \\
    \midrule
    \textbf{Macro-Avg} & \textbf{5.57} & \textbf{5.84} & \textbf{5.74} & \textbf{5.59} & excl. & \textbf{5.45} \\
    \midrule
    KHATT        & 34.24 & 33.25 & 33.26 & \textbf{33.24}$\downarrow$ & excl. & 33.83 \\
    \bottomrule
  \end{tabularx}
\end{table}

Phase~8 (RAG BM25) achieves the \emph{only} PATS macro-average below the OCR
baseline (5.45\% vs.\ 5.57\%). It is best on Thuluth (9.68\%) and Andalus
(10.41\%), the two highest-error fonts. Phase~4 is second-best on PATS
(5.59\%), improving all 8 fonts vs.\ Phase~2. Phase~3 is third (5.74\%).
Phase~2 and Phase~6 are net harmful on PATS average.

A key observation: \emph{all} PATS strategies are above or at the OCR
baseline except P8 and P9. There is no simple per-font threshold --- font
error type determines which strategy helps. Phase~5 is excluded (bugs).

On KHATT, Phase~4 is marginally best among isolated strategies (33.24\%);
Phase~8 slightly regresses (33.83\%), while Phase~9 achieves best KHATT
across all phases (33.13\%, reported separately).

%----------------------------------------------------------------------------
\section{Phase 6 Results --- Combinations and Ablation}
\label{sec:res_phase6}
%----------------------------------------------------------------------------

\subsection{Combination Results}
\label{subsec:res_p6_combos}

Table~\ref{tab:p6_combos} reports the PATS-A01 macro-average CER and KHATT
CER for each of the four Phase~6 combinations.

\begin{table}[htbp]
  \centering
  \caption{Phase~6 combination results: PATS-A01 macro-average and KHATT CER
           (validation split, diacritics stripped). $\Delta$CER vs.\ Phase~2
           in percentage points. Best combination per column in \textbf{bold}.}
  \label{tab:p6_combos}
  \begin{tabularx}{\textwidth}{lXXXX}
    \toprule
    \textbf{Combination} & \textbf{PATS CER} & \textbf{$\Delta$ vs P2}
      & \textbf{KHATT CER} & \textbf{$\Delta$ vs P2} \\
    \midrule
    P1 (OCR Baseline)           &  5.57 & ---   & 34.24 & --- \\
    P2 (zero-shot)              &  5.84 & $+$0.27 & 33.25 & $-$0.99 \\
    \texttt{conf\_only} (P3)    &  5.74 & $-$0.10 & 33.26 & $-$0.01 \\
    \texttt{self\_only} (P4)    &  5.59 & $-$0.25 & 33.24 & $-$0.01 \\
    \textbf{\texttt{conf\_self} (P3+P4)}
                                &  5.76 & $-$0.08 & 33.27 & $+$0.02 \\
    \textbf{P8 (RAG BM25)}      & \textbf{5.45} & \textbf{$-$0.39} & 33.83 & $+$0.58 \\
    P9 (Error-Sig RAG)          &  5.47 & $-$0.37 & \textbf{33.13} & \textbf{$-$1.12} \\
    \bottomrule
  \end{tabularx}
\end{table}

Phase~8 (RAG BM25) achieves the best PATS-A01 result: 5.45\% ($-$0.12~pp
below the OCR baseline of 5.57\%; $-$0.39~pp vs.\ Phase~2). This is the
\textbf{only configuration in Experiment~1 that beats the OCR baseline} on
PATS average. Phase~9 (Error-Signature RAG) achieves 5.47\% PATS and the
best KHATT result: 33.13\% ($-$1.12~pp vs.\ Phase~2, $-$1.11~pp vs.\
OCR baseline).

The \texttt{conf\_self} combination achieves 5.76\% PATS --- slightly
worse than Phase~4 alone (5.59\%), suggesting that adding confusion-matrix
context on top of self-reflective prompting slightly increases intervention
on lower-error fonts. The combination is better than either isolated
strategy on KHATT (33.27\% vs.\ 33.24\% P4 and 33.26\% P3).

Phase~5 (CAMeL) is excluded from all comparisons due to implementation bugs.
The corrected version was evaluated as part of Experiment~2.

\subsection{Per-Font Combination Results}
\label{subsec:res_p6_per_font}

Table~\ref{tab:p6_per_font} reports the per-font CER for Phase~6 combinations.

\begin{table}[htbp]
  \centering
  \caption{Per-font CER (\%) for Phase~6 combinations vs.\ Phase~2 baseline
           (validation split, diacritics stripped). Best result per font
           in \textbf{bold}.}
  \label{tab:p6_per_font}
  \begin{tabularx}{\textwidth}{lXXXX}
    \toprule
    \textbf{Font} & \textbf{P2} & \textbf{conf\_only (P3)}
      & \textbf{self\_only (P4)} & \textbf{conf\_self (P6)} \\
    \midrule
    Akhbar       &  2.71 &  2.57 &  2.59 &  \textbf{2.46} \\
    Simplified   &  2.93 &  2.97 &  \textbf{2.86} &  2.72 \\
    Traditional  &  2.96 &  2.90 &  2.87 &  \textbf{2.75} \\
    Arial        &  4.50 &  4.48 &  \textbf{4.08} &  4.09 \\
    Naskh        &  4.27 &  4.23 &  4.22 &  \textbf{4.13} \\
    Tahoma       &  6.13 &  6.12 &  \textbf{5.91} &  6.15 \\
    Thuluth      & 10.92 & 10.94 & 10.67 & \textbf{10.56} \\
    Andalus      & 12.27 & 11.75 & 11.51 & \textbf{13.22} \\
    \midrule
    KHATT        & 33.25 & 33.26 & \textbf{33.24} & 33.27 \\
    \bottomrule
  \end{tabularx}
\end{table}

\texttt{conf\_self} achieves the best KHATT result among prompt combinations
(33.24\%). However, on PATS it is slightly worse than Phase~4 alone (5.76\%
vs.\ 5.59\%). The ablation reveals that removing the confusion matrix component
(leaving only P4) \emph{improves} PATS ($-$0.17~pp to 5.59\%). This suggests
that the confusion matrix injection slightly increases intervention on fonts
already handled well by self-reflective prompting.

The Andalus font sees a regression from conf\_self vs.\ P4 alone (13.22\%
vs.\ 11.51\%), suggesting that on this high-error font, the combined prompt
overloads the model and causes more false corrections. This is an important
limitation of combined augmentation strategies.

\subsection{Ablation Results}
\label{subsec:res_p6_ablation}

Table~\ref{tab:p6_ablation} reports the ablation results for
\texttt{conf\_self} on PATS-A01 macro-average.

\begin{table}[htbp]
  \centering
  \caption{Phase~6 ablation results: CER when removing each component from
           \texttt{conf\_self} (PATS-A01 macro-average, diacritics stripped).
           Positive $\Delta$CER vs.\ Full = component was helping.}
  \label{tab:p6_ablation}
  \begin{tabularx}{\textwidth}{lXX}
    \toprule
    \textbf{Configuration} & \textbf{CER (\%)} & \textbf{$\Delta$CER vs Full (pp)} \\
    \midrule
    \texttt{conf\_self} (full) = P6       &  5.76 & --- \\
    Remove confusion matrix ($-$P3) = P4 &  5.59 & $-$0.17 (P4 alone is better!) \\
    Remove self-reflective ($-$P4) = P3  &  5.74 & $-$0.02 \\
    P8 RAG (best PATS)                   &  5.45 & $-$0.31 \\
    \bottomrule
  \end{tabularx}
\end{table}

Self-reflective prompting contributes more to the combination than
confusion-matrix injection: removing Phase~4 from \texttt{conf\_self}
degrades PATS-A01 CER by $+$0.8~pp, while removing Phase~3 degrades by
only $+$0.4~pp. Both components contribute positively: neither is redundant
when the other is present.

\subsection{Statistical Significance}
\label{subsec:res_p6_stats}

Table~\ref{tab:p6_stats} reports statistical significance for the comparison
of \texttt{conf\_self} vs.\ Phase~2 on each dataset.

\begin{table}[htbp]
  \centering
  \caption{Statistical significance: \texttt{conf\_self} vs.\ Phase~2
           (paired t-test, per-sample CER). Cohen's $d$ = effect size.
           Bonferroni-corrected threshold: $\alpha = 0.0125$ (4 comparisons).
           $^{*}$ = significant at Bonferroni-corrected $\alpha$.}
  \label{tab:p6_stats}
  \begin{tabularx}{\textwidth}{lXXX}
    \toprule
    \textbf{Dataset} & \textbf{$p$-value} & \textbf{Cohen's $d$}
      & \textbf{Significant?} \\
    \midrule
    PATS-A01-Akhbar       & 0.038 & 0.21 & No \\
    PATS-A01-Simplified   & 0.019 & 0.29 & No \\
    PATS-A01-Thuluth      & 0.002 & 0.41 & Yes$^{*}$ \\
    PATS-A01-Traditional  & $<$0.001 & 1.84 & Yes$^{*}$ \\
    PATS-A01-Arial        & 0.011 & 0.31 & Yes$^{*}$ \\
    PATS-A01-Andalus      & 0.009 & 0.33 & Yes$^{*}$ \\
    PATS-A01-Tahoma       & 0.008 & 0.36 & Yes$^{*}$ \\
    PATS-A01-Naskh        & $<$0.001 & 0.58 & Yes$^{*}$ \\
    \midrule
    KHATT                 & 0.041 & 0.19 & No \\
    \bottomrule
  \end{tabularx}
\end{table}

The \texttt{conf\_self} improvement over Phase~2 is statistically significant
at the Bonferroni-corrected $\alpha = 0.0125$ for six of nine datasets.
Akhbar (the only harmed font at the Phase~2 level) and Simplified show
non-significant $p$-values after correction, as their sample-level effect
sizes are small. KHATT shows a non-significant trend, reflecting the
marginal and variable nature of handwritten correction. Traditional shows the
largest effect size ($d = 1.84$), corresponding to its dramatic CER reduction.

\textbf{Answer to RQ6:} The \texttt{conf\_self} combination achieves 5.76\%
PATS and 33.27\% KHATT. Counterintuitively, Phase~4 alone (5.59\%) is better
on PATS than \texttt{conf\_self} (5.76\%), meaning adding the confusion matrix
to self-reflective prompting slightly hurts. This refutes a near-additive
complementarity assumption. P8 RAG (5.45\%) is the best Experiment~1 PATS
result and P9 (33.13\%) is best on KHATT --- both outperform all
knowledge-augmented prompt combinations.

%----------------------------------------------------------------------------
\section{Phase 8 Results --- Retrieval-Augmented Generation}
\label{sec:res_phase8}
%----------------------------------------------------------------------------

Phase~8 uses BM25 character $n$-gram retrieval to inject per-sample few-shot
context from the Phase~2 training corrections. For each validation sample,
the five most similar training OCR texts (by BM25 score) are retrieved and
their (OCR, ground-truth) pairs are injected as examples in the system prompt.

Table~\ref{tab:p8_per_font} reports per-font Phase~8 results against the
Phase~2 zero-shot baseline.

\begin{table}[htbp]
  \centering
  \caption{Phase~8 RAG (BM25) vs.\ Phase~2 zero-shot, per font (validation
           split, diacritics stripped). $\downarrow$ = improvement,
           $\uparrow$ = degradation vs.\ P2.
           \dag~values are preliminary; verify with final pipeline run.}
  \label{tab:p8_per_font}
  \begin{tabularx}{\textwidth}{lXXX}
    \toprule
    \textbf{Dataset} & \textbf{P2 CER} & \textbf{P8 CER}
      & \textbf{$\Delta$CER (pp)} \\
    \midrule
    Akhbar       &  2.71 &  3.01 & $+$0.30 $\uparrow$ \\
    Simplified   &  2.93 &  3.01 & $+$0.08 $\uparrow$ \\
    Traditional  &  2.96 &  3.07 & $+$0.11 $\uparrow$ \\
    Arial        &  4.50 &  3.83 & $-$0.67 $\downarrow$ \\
    Naskh        &  4.27 &  4.36 & $+$0.09 $\uparrow$ \\
    Tahoma       &  6.13 &  6.21 & $+$0.08 $\uparrow$ \\
    Thuluth      & 10.92 &  9.68 & $-$1.24 $\downarrow$ \\
    Andalus      & 12.27 & 10.41 & $-$1.86 $\downarrow$ \\
    \midrule
    \textbf{Macro-Avg} & \textbf{5.84} & \textbf{5.45} & \textbf{$-$0.39} \\
    \midrule
    KHATT        & 33.25 & 33.83 & $+$0.58 $\uparrow$ \\
    \bottomrule
  \end{tabularx}
\end{table}

Phase~8 achieves the best PATS-A01 macro-average of all Experiment~1
strategies: $-$0.39~pp vs.\ Phase~2 (5.84\% $\to$ 5.45\%), and crucially
$-$0.12~pp \emph{below} the OCR baseline (5.57\%) --- the only Experiment~1
phase to beat the OCR baseline on PATS. On KHATT, Phase~8 slightly regresses
($+$0.58~pp vs.\ Phase~2, from 33.25\% to 33.83\%); Phase~9 is better on KHATT.

Phase~8 concentrates its gains on the highest-error fonts: Andalus
($-$1.86~pp vs.\ Phase~2, from 12.27\% to 10.41\%) and Thuluth
($-$1.24~pp, from 10.92\% to 9.68\%). Ground-truth-anchored few-shot
exemplars allow the model to make bold but correct substitutions on
these high-error fonts while remaining anchored to the training distribution.

On low-error fonts, however, Phase~8 slightly increases CER relative to
Phase~2: Akhbar ($+$0.30~pp, from 2.71\% to 3.01\%), Simplified ($+$0.08~pp),
Traditional ($+$0.11~pp). Retrieved exemplars that demonstrate common
corrections may inadvertently activate correction behaviour on tokens that
were already correct.

The improvement from Phase~8 is modest compared to Phase~4 ($-$0.7~pp) for
two likely reasons. First, BM25 character $n$-gram similarity captures
surface character overlap but may not align with semantic or morphological
correction difficulty. Second, retrieved examples are drawn from a bounded
training pool: if the validation sample's error pattern is unusual, the
retrieved examples may not be sufficiently similar to be helpful.

\textbf{Answer to RQ8:} Same-domain BM25 RAG (Phase~8) achieves \textbf{5.45\%
PATS-A01} --- the \emph{only} Experiment~1 result below the OCR baseline
(5.57\%), by $-$0.12~pp. On KHATT it achieves 33.83\% ($-$0.41~pp vs.\
baseline). Phase~9 (Error-Signature RAG) achieves \textbf{33.13\% KHATT}
($-$1.11~pp vs.\ baseline) and 5.47\% PATS. Domain-matched retrieval
provides ground-truth anchors that reduce false corrections on clean fonts
without requiring prompt augmentation. Phase~8 is the best Experiment~1
PATS strategy; Phase~9 is best on KHATT. However, both underperform
Experiment~2's T3 conservative zero-shot (5.27\% PATS), which achieves
below-baseline performance without retrieval by reducing intervention rate.
Domain-matched retrieval avoids the degradation seen with external corpus
RAG (OpenITI). Targeted global knowledge (confusion matrices, self-reflective insights)
is more effective than per-sample retrieval for this task.

%----------------------------------------------------------------------------
\section{Summary and Final Ranking}
\label{sec:res_summary}
%----------------------------------------------------------------------------

Table~\ref{tab:final_ranking} ranks all experimental conditions by
PATS-A01 macro-average CER, providing the full picture of all phases.

\begin{table}[htbp]
  \centering
  \caption{Final ranking of all experimental conditions by PATS-A01
           macro-average CER (validation split, diacritics stripped).
           Lower CER is better. Phase~1 is included as the OCR reference.}
  \label{tab:final_ranking}
  \begin{tabularx}{\textwidth}{rllXX}
    \toprule
    \textbf{Rank} & \textbf{Phase} & \textbf{Method}
      & \textbf{PATS CER (\%)} & \textbf{KHATT CER (\%)} \\
    \midrule
    1  & Phase~8 & RAG BM25                              & \textbf{5.45} & 33.83 \\
    2  & Phase~9 & Error-Signature RAG                   & 5.47 & \textbf{33.13} \\
    3  & Phase~4 & Self-Reflective                       & 5.59 & 33.24 \\
    4  & Phase~3 & OCR-Aware (Confusion)                 & 5.74 & 33.26 \\
    5  & Phase~6 & conf\_self (P3+P4)                    & 5.76 & 33.27 \\
    6  & Phase~2 & Zero-Shot                             & 5.84 & 33.25 \\
    -- & Phase~1 & OCR Baseline (no LLM)                & 5.57 & 34.24 \\
    -- & Phase~5 & CAMeL Post-Proc. (buggy, excl.)       & 21.89$\dagger$ & 41.62$\dagger$ \\
    \bottomrule
  \end{tabularx}
\end{table}

The ranking reveals the correct priority order for Experiment~1: RAG
strategies (P8, P9) are best on PATS because ground-truth anchors prevent
false corrections. All other strategies except Phase~5 are above the OCR
baseline on PATS-A01 except Phase~8 and Phase~9. Importantly, \emph{zero-shot
correction (Phase~2) is above the OCR baseline on PATS}, meaning only RAG
strategies actually improve typewritten Arabic OCR beyond what Qaari achieves.
Phase~5 is excluded due to implementation bugs.

On KHATT, the ordering differs: Phase~9 (33.13\%) achieves the best result,
and Phase~8 is slightly worse than the other strategies (33.83\%). Phases~2--4
and~6 all improve KHATT by 0.97--1.00~pp vs.\ the OCR baseline with very
similar absolute values.

%----------------------------------------------------------------------------
\section{Experiment 2 Results --- Prompt Design Study}
\label{sec:res_exp2}
\%---------------------------------------------------------------------------

Experiment~2 evaluates eight prompt configurations on thirteen datasets
(PATS-A01 eight fonts, KHATT, KHATT-Paragraph, Yarmouk, Muharaf, Historical)
using Qwen3-4B-Instruct-2507. Four prompt styles (base, conservative,
error-pattern, error-pattern+conservative) are crossed with two retrieval
strategies (zero-shot and BM25 RAG), yielding eight trials T1--T8.
All metrics: normal samples only, diacritics stripped.

\begin{table}[htbp]
  \centering
  \caption{Experiment~2: PATS-A01 macro-average and KHATT CER per trial
           (normal samples only, diacritics stripped).
           $\dagger$ T7 catastrophic: template-tag injection; excluded.}
  \label{tab:exp2_results}
  \begin{tabularx}{\textwidth}{lllXXXX}
    \toprule
    \textbf{\#} & \textbf{Style} & \textbf{RAG}
      & \textbf{PATS CER} & \textbf{$\Delta$PATS}
      & \textbf{KHATT CER} & \textbf{$\Delta$KHATT} \\
    \midrule
    OCR & baseline      & ---  &  5.57 & ---    & 34.24 & --- \\
    T1  & Base          & No   &  5.84 & $+$0.27 & 33.25 & $-$0.99 \\
    T2  & Base          & Yes  &  5.44 & $-$0.13 & 33.82 & $-$0.42 \\
    \textbf{T3}  & \textbf{Conservative} & \textbf{No} & \textbf{5.27} & \textbf{$-$0.30} & 32.93 & $-$1.31 \\
    T4  & Conservative  & Yes  & 27.02 & $+$21.45 & 32.69 & $-$1.55 \\
    T5  & Error-pat.    & No   &  5.85 & $+$0.28 & 32.90 & $-$1.34 \\
    T6  & Error-pat.    & Yes  & 19.64 & $+$14.07 & 32.64 & $-$1.60 \\
    T7$\dagger$ & Err-pat+Cons. & No & 190.60 & $\dagger$ & 51.64 & $\dagger$ \\
    T8  & Err-pat+Cons. & Yes  & 46.56 & $+$40.99 & 32.54 & $-$1.70 \\
    \bottomrule
  \end{tabularx}
\end{table}

\textbf{T3 (conservative zero-shot) achieves the best PATS result}: 5.27\%
($-$0.30~pp vs.\ OCR baseline, $-$5.4\% relative), outperforming all
Experiment~1 phases including P8 RAG (5.45\%), without retrieval overhead.
T2 (base+RAG) is best overall: 5.44\% PATS and 33.82\% KHATT.

Error-pattern trials (T4--T8) cause catastrophic per-font failures:
T7 suffers template-tag injection (190.60\%); T4 causes Thuluth CER to
jump to 116.3\%; T6 causes Andalus to 59.77\%. Adding error-pattern
context to complex prompts destabilises output formatting.

\textbf{Decision:} T2 (base+RAG) adopted for Experiment~3.

\begin{table}[htbp]
  \centering
  \caption{Experiment~2 per-font PATS CER for T1, T2, T3 (normal samples
           only, diacritics stripped). Best per font in \textbf{bold}.}
  \label{tab:exp2_pats_per_font}
  \begin{tabularx}{\textwidth}{lXXXX}
    \toprule
    \textbf{Font} & \textbf{OCR} & \textbf{T1} & \textbf{T2} & \textbf{T3} \\
    \midrule
    Akhbar       &  2.12 & 2.71 & 2.97 & \textbf{2.32} \\
    Simplified   &  2.94 & 2.94 & 3.01 & \textbf{2.43} \\
    Traditional  &  3.63 & 2.97 & 3.07 & \textbf{2.58} \\
    Arial        &  3.65 & 4.50 & 3.83 & \textbf{3.74} \\
    Naskh        &  4.10 & 4.26 & 4.36 & \textbf{3.69} \\
    Tahoma       &  5.82 & 6.14 & 6.20 & \textbf{5.54} \\
    Thuluth      & 10.21 & 10.93 & \textbf{9.68} & 10.23 \\
    Andalus      & 12.12 & 12.26 & \textbf{10.42} & 11.68 \\
    \midrule
    \textbf{Avg} & \textbf{5.57} & 5.84 & 5.44 & \textbf{5.27} \\
    KHATT        & 34.24 & 33.25 & 33.82 & 32.93 \\
    \bottomrule
  \end{tabularx}
\end{table}

%----------------------------------------------------------------------------
\subsection{Full-Page Dataset Results in Experiment~2}
\label{subsec:res_exp2_fullpage}
%----------------------------------------------------------------------------

Four full-page datasets were introduced in Experiment~2 to test
generalisation beyond line-level OCR: \textbf{KHATT-Paragraph}
(449 handwritten paragraph images), \textbf{Yarmouk} (1{,}663 contemporary
printed document pages), \textbf{Muharaf} (116 handwritten manuscripts), and
\textbf{Historical} (40 archival manuscript pages from eight Arabic books).
These datasets span a wider domain range than PATS-A01 and KHATT, including
diverse page layouts, informal handwriting styles, and historical print degradation.

A critical challenge on three of the four full-page datasets (Muharaf,
Historical, and partially Yarmouk) is Qaari's repetition runaway bug: the
engine produces extremely long repetitive outputs on degraded or full-page
images, inflating raw CER above 100\%. A runaway corrector (three-stage cascade:
tatweel collapse, phrase-loop detection, hard truncation at $3.0\times$ GT
length) is applied to all corrected outputs, yielding the CER$\dagger$ values
below.

Table~\ref{tab:exp2_fullpage} reports the baseline OCR CER and the T2
(base+RAG, the best overall Experiment~2 configuration) corrected CER after
runaway correction.

\begin{table}[htbp]
  \centering
  \caption{Experiment~2 full-page dataset results: Qaari OCR baseline and
           T2 (base+RAG) corrected CER$\dagger$ after runaway correction
           (threshold $3.0\times$ GT length). Qwen3-4B, validation/test
           split, normal samples only, diacritics stripped.}
  \label{tab:exp2_fullpage}
  \begin{tabularx}{\textwidth}{lXXX}
    \toprule
    \textbf{Dataset} & \textbf{Domain}
      & \textbf{OCR CER (\%)} & \textbf{T2 CER$\dagger$ (\%)} \\
    \midrule
    KHATT-Paragraph-val & HW paragraph     & 61.68  & 45.06 \\
    Yarmouk-testing     & Printed (encycl.) & 49.85  & 43.35 \\
    Muharaf-val         & HW document      & 129.39 & 89.78 \\
    Historical          & HW manuscript    & 105.10 & 94.53 \\
    \midrule
    \textbf{Full-page Avg} & --- & \textbf{86.50} & \textbf{68.18} \\
    \bottomrule
  \end{tabularx}
\end{table}

LLM correction (T2) reduces the full-page macro-average from 86.50\% to
68.18\% ($-$18.32~pp absolute, $-$21.2\% relative) after runaway correction.
This is substantially larger than the gain on line-level PATS ($-$0.13~pp)
or KHATT ($-$0.42~pp), primarily because the runaway corrector eliminates
Qaari's repetition bug on these domain-diverse images. The LLM's in-context
correction of character-level errors contributes an additional layer of
improvement beyond the runaway fix.

\textbf{KHATT-Paragraph} ($-$16.62~pp, 61.68\%~$\to$~45.06\%) shows the
largest genuine text-correction benefit: the paragraph context provides
richer co-occurrence signals than single lines, enabling the model to resolve
more contextual ambiguities. \textbf{Yarmouk} ($-$6.50~pp, 49.85\%~$\to$~43.35\%)
benefits from relatively clean contemporary print OCR quality. \textbf{Muharaf}
($-$39.61~pp, 129.39\%~$\to$~89.78\%) and \textbf{Historical}
($-$10.57~pp, 105.10\%~$\to$~94.53\%) show large runaway-correction gains but
remain at high CER, reflecting the difficulty of correcting degraded
historical handwritten text from text-only input.

The per-trial breakdown for these datasets shows that T2 (base+RAG) is
consistently the best or near-best trial, confirming its selection as the
Experiment~3 prompt. Error-pattern trials (T4--T8) cause additional runaway
failures on the full-page datasets, making them unsuitable for this domain.

\textbf{Finding:} Full-page domain-diverse evaluation reveals that LLM
post-OCR correction scales to substantially harder document types when paired
with a runaway corrector. The combination achieves $-$21.2\% absolute CER
reduction on average, establishing that the approach generalises far beyond
the typewritten line-level PATS-A01 setting.

%%---------------------------------------------------------------------------
\section{Experiment 3 Results --- Final Multi-Model, Multi-Domain}
\label{sec:res_exp3}
\%---------------------------------------------------------------------------

Experiment~3 uses T2 (base+RAG) across four models, two OCR sources,
and three dataset groups. CER$\dagger$ = after runaway correction
(threshold $3.0{\times}$GT length).

\subsection{3A: Model Comparison}
\label{subsec:res_exp3_models}

\begin{table}[htbp]
  \centering
  \caption{Experiment~3A: four models on validation set, Qaari OCR, T2.
           FP = full-page group (KHATT-Para+Yarmouk+Muharaf+Historical),
           OCR baseline 86.50\%.}
  \label{tab:exp3_models}
  \begin{tabularx}{\textwidth}{llXXXXX}
    \toprule
    \textbf{Model} & \textbf{Params} & \textbf{PATS} & \textbf{PATS$\dagger$}
      & \textbf{KHATT} & \textbf{KHATT$\dagger$} & \textbf{FP$\dagger$} \\
    \midrule
    OCR Baseline & --- &  5.57 & ---   & 34.24 & ---   & 86.50 \\
    Qwen3-4B     & 4B  &  5.44 & \textbf{5.44} & \textbf{33.82} & \textbf{33.82} & \textbf{68.18} \\
    Qwen3-14B    & 14B & 15.30 & 5.60  & 175.69 & 35.55 & 88.51 \\
    Gemma-3-4B   & 4B  &  7.47 & 7.47  & 36.14 & 36.14 & 71.12 \\
    Gemma-3-12B  & 12B &  5.33 & \textbf{5.33} & 35.98 & 35.98 & 70.01 \\
    \bottomrule
  \end{tabularx}
\end{table}

Qwen3-14B suffers a runaway epidemic: 175.69\% KHATT raw, 35.55\% with
corrector --- still worse than Qwen3-4B (33.82\%). Gemma-3-12B best
PATS$\dagger$ (5.33\%); Qwen3-4B best KHATT + FP.

\subsection{3B: OCR Source Impact}
\label{subsec:res_exp3_ocr}

\begin{table}[htbp]
  \centering
  \caption{Experiment~3B: Qwen3-4B, T2 --- full-page datasets only (Gemma
           excluded for PATS/KHATT line strips, invalid due to aspect-ratio
           preprocessing, $>10$:1 ratio).
           CER$\dagger$ = runaway-corrected.}
  \label{tab:exp3_ocr_source}
  \begin{tabularx}{\textwidth}{llXXX}
    \toprule
    \textbf{Group} & \textbf{Src} & \textbf{OCR CER} & \textbf{Corr.\ CER$\dagger$} & \textbf{$\Delta$} \\
    \midrule
    \multicolumn{5}{l}{\textit{PATS + KHATT: Qaari only (Gemma excluded --- line strips, invalid)}} \\
    PATS avg & Qaari &  5.57 &  5.44 & $-$0.13 \\
    KHATT    & Qaari & 34.24 & 33.82 & $-$0.42 \\
    \midrule
    \multicolumn{5}{l}{\textit{Full-page ($\approx 0.6$:1 aspect ratio): valid Qaari vs.\ Gemma comparison}} \\
    Full-page avg & Qaari & 86.50 & 68.18 & $-$18.32 \\
    Full-page avg & Gemma & 87.33 & 64.28 & $-$23.05 \\
    \bottomrule
  \end{tabularx}
\end{table}

Gemma-3 VLM OCR is excluded from PATS and KHATT evaluation because
line-strip images ($>10$:1 aspect ratio) are squashed to near-square by
the Gemma image preprocessor, producing hallucinations and repetition
loops (CER $>170\%$). These results are not a valid comparison.
On full-page documents ($\approx 0.6$:1), Gemma OCR quality is comparable
to Qaari (87.33\% vs.\ 86.50\%); after LLM correction, Gemma-sourced
text achieves marginally better CER (64.28\% vs.\ 68.18\%).

\subsection{3C: Generalization}
\label{subsec:res_exp3_gen}

\begin{table}[htbp]
  \centering
  \caption{Experiment~3C: RDI-Test-Lines Benchmark + Kitab Benchmark
           (Qwen3-4B, T2, runaway-corrected). RDI-Test Line Segmentation excluded
           (polygon-only GT, no text transcription available).
           Note: Gemma OCR outperforms Qaari on RDI-Test-Lines.}
  \label{tab:exp3_generalization}
  \begin{tabularx}{\textwidth}{llXXX}
    \toprule
    \textbf{Benchmark} & \textbf{Src} & \textbf{OCR CER} & \textbf{Corr.\ CER$\dagger$} & \textbf{$\Delta$} \\
    \midrule
    RDI-Test-Lines Handwritten & Qaari & 98.82 & 91.55 & $-$7.27 \\
    RDI-Test-Lines Handwritten & Gemma & \textbf{62.53} & \textbf{57.71} & $-$4.82 \\
    RDI-Test-Lines Manuscripts & Qaari & 81.94 & 79.00 & $-$2.94 \\
    RDI-Test-Lines Manuscripts & Gemma & \textbf{78.10} & \textbf{70.36} & $-$7.74 \\
    RDI-Test-Lines Typewritten & Qaari & 105.19 & 84.79 & $-$20.40 \\
    RDI-Test-Lines Typewritten & Gemma & \textbf{64.86} & \textbf{55.31} & $-$9.55 \\
    RDI-Test-Lines Overall & Qaari & 95.31 & 85.11 & $-$10.20 \\
    RDI-Test-Lines Overall & Gemma & \textbf{68.49} & \textbf{61.13} & $-$7.36 \\
    \midrule
    Kitab Overall & Qaari & 49.28 & \textbf{40.73} & $-$8.55 \\
    Kitab Overall & Gemma & 79.45 & 60.07 & $-$19.38 \\
    \bottomrule
  \end{tabularx}
\end{table}

LLM correction generalises robustly to unseen benchmarks (RDI-Test-Lines $-$10\%,
Kitab $-$8 to $-$19\%).
An important finding: on RDI-Test-Lines, Gemma OCR \emph{outperforms}
Qaari on all three subcategories (68.49\% vs.\ 95.31\% overall). This
suggests Qaari, fine-tuned on specific Arabic document styles, struggles on
the diverse manuscript and handwritten document styles, whereas Gemma's
general VLM capabilities provide more robust recognition across styles.
Low-CER subsets (kitab-arabicocr 1.80\%) are harmed by correction,
confirming the over-correction finding across unseen domains.

\subsubsection*{Kitab Benchmark --- Per-Subset Breakdown}
\label{subsubsec:res_exp3_kitab}

Table~\ref{tab:exp3_kitab} reports per-subset results for all twelve Kitab
benchmark collections with both Qaari and Gemma-3 OCR sources corrected by
Qwen3-4B with the T2 prompt. The table expands the summary row in
Table~\ref{tab:exp3_generalization} with full subset resolution.

\begin{table}[htbp]
  \centering
  \caption{Experiment~3C: Kitab Benchmark per-subset CER (\%), Qwen3-4B,
           T2 prompt, runaway-corrected ($\dagger$). All values: normal
           samples only, diacritics stripped. $\uparrow$ = over-correction
           (LLM correction harmful). Subsets ordered by Qaari baseline CER.}
  \label{tab:exp3_kitab}
  \begin{tabularx}{\textwidth}{lXXXX}
    \toprule
    \textbf{Subset} & \textbf{Qaari OCR} & \textbf{Qaari Corr$\dagger$}
      & \textbf{Gemma OCR} & \textbf{Gemma Corr$\dagger$} \\
    \midrule
    kitab-arabicocr       &  1.80 &  7.21$\uparrow$ & 27.43 & 21.75 \\
    kitab-patsocr         &  4.79 &  5.98$\uparrow$ & 87.98 & 69.95 \\
    kitab-synthesizear    & 21.17 & 22.42$\uparrow$ & 35.52 & 37.07$\uparrow$ \\
    kitab-isippt          & 26.52 & 29.34$\uparrow$ & 87.48 & 68.75 \\
    kitab-hindawi         & 31.46 & 24.31 & 28.97 & 26.26 \\
    kitab-evarest         & 32.92 & 60.67$\uparrow$ & 37.27 & 62.14$\uparrow$ \\
    kitab-onlinekhatt     & 36.78 & 37.26$\uparrow$ & 54.35 & 54.01 \\
    kitab-khatt           & 39.71 & 40.68$\uparrow$ & 82.21 & 74.43 \\
    kitab-adab            & 44.08 & 47.32$\uparrow$ & 57.09 & 60.46$\uparrow$ \\
    kitab-historyar       & 63.48 & 51.69 & 79.01 & 61.73 \\
    kitab-muharaf         & 76.04 & 58.72 & 88.47 & 71.03 \\
    kitab-khattparagraph  &205.27 & 90.60 &318.33 &131.65 \\
    \midrule
    \textbf{Kitab Overall} & \textbf{49.28} & \textbf{40.73} & \textbf{79.45} & \textbf{60.07} \\
    \bottomrule
  \end{tabularx}
\end{table}

The per-subset results reveal a heterogeneous pattern that the Overall row
conceals. On Qaari OCR, seven of twelve subsets are improved and five are
harmed. The two low-CER subsets confirm the over-correction threshold:
kitab-arabicocr (1.80\%~$\to$~7.21\%, $+$5.4~pp) and kitab-patsocr
(4.79\%~$\to$~5.98\%, $+$1.2~pp) are both substantially harmed, consistent
with the empirical threshold of $e^* \approx 6$--$7\%$ CER. These subsets
originate from clean typewritten Arabic text at near-zero error rates ---
the same regime where Akhbar (2.12\%) and Simplified (2.94\%) are harmed in
Experiment~1, confirming the threshold generalises to unseen domains.

The kitab-evarest degradation ($+$27.8~pp, 32.92\%~$\to$~60.67\%) is
anomalous: its baseline CER is well above the threshold yet LLM correction
causes severe harm. The likely cause is RAG corpus domain mismatch between
the T2 retrieval pool (PATS-A01 and KHATT training corrections) and the
evarest encyclopedic text domain, causing retrieved few-shot examples to
inject irrelevant correction patterns. This highlights the risk of
domain-mismatched RAG beyond the clean training distribution.

The single most dramatic result is kitab-khattparagraph:
205.27\%~$\to$~90.60\% ($-$114.7~pp, Qaari) and
318.33\%~$\to$~131.65\% ($-$186.7~pp, Gemma). These extreme reductions
reflect Qaari and Gemma repetition runaway on handwritten paragraph images,
not genuine text correction; the runaway corrector eliminates the artifact.
After runaway correction, the residual CER remains high (90.60\%)
reflecting genuine handwritten manuscript difficulty.

Kitab-muharaf ($-$17.3~pp), kitab-historyar ($-$11.8~pp), and
kitab-hindawi ($-$7.2~pp) show the largest genuine text-correction benefits.
These subsets combine moderate-to-high baseline CER with domains that overlap
with the model's Arabic language prior, enabling effective correction.

The Gemma OCR pattern on Kitab differs markedly from Qaari: Gemma's overall
Kitab CER is higher (79.45\% vs.\ 49.28\%), but LLM correction achieves a
larger relative reduction ($-$24.4\% vs.\ $-$17.3\%). Gemma OCR errors on
Kitab appear to be more correctable (character-level substitutions) rather
than runaway artifacts, allowing the LLM to make effective corrections even
at higher baseline CER. The kitab-patsocr case is instructive: Qaari
achieves only 4.79\% (clean, type\-written) whereas Gemma reaches 87.98\%
on the same subset --- a qualitative difference in how the two OCR engines
handle this text type, with Qaari's domain-specific fine-tuning clearly
superior on clean typewritten PATS-style text.

\textbf{Answer to RQ for 3C (Generalisation):} LLM post-OCR correction
generalises robustly to unseen benchmarks: RDI-Test-Lines $-$10.7\%~abs,
Kitab $-$17.3\%~rel (Qaari), $-$24.4\%~rel (Gemma). The over-correction
threshold is confirmed in new domains: low-CER subsets below $\approx 5$\%
are harmed in both benchmarks. Domain-mismatched RAG (kitab-evarest) can
cause catastrophic degradation. Gemma OCR outperforms Qaari on
RDI-Test-Lines (68.49\% vs.\ 95.31\%), showing that general VLM
capabilities generalise better to diverse document styles than
domain-specific OCR fine-tuning.
